\section{TỔNG KẾT}

Trong bối cảnh chuyển đổi số ngày càng lan rộng và đóng vai trò then chốt trong lĩnh vực y tế, đề tài \textbf{\textit{“Xây dựng hệ thống chăm sóc sức khỏe điện tử với thiết bị theo dõi sức khỏe thông minh”}} đã được triển khai với kỳ vọng góp phần đổi mới phương thức theo dõi và quản lý sức khỏe cá nhân, đồng thời hỗ trợ hiệu quả cho các cơ sở y tế trong công tác giám sát cộng đồng. Qua quá trình thực hiện, hệ thống đã tích hợp thành công giữa phần cứng (thiết bị cân thông minh) và phần mềm (nền tảng phân tích và hiển thị dữ liệu sức khỏe), tạo nên một giải pháp có tính ứng dụng cao và phù hợp với thực tiễn.

\subsection{Đánh giá kết quả thực hiện}

Về mặt kỹ thuật, hệ thống đã đạt được các yêu cầu đặt ra trong đề tài. Cụ thể:

\begin{itemize}
  \item \textbf{Khả năng thu thập và xử lý dữ liệu}: Hệ thống đã kết nối thành công với các thiết bị cân sức khỏe thông minh thông qua Bluetooth Low Energy, trích xuất chính xác các thông số như cân nặng, tỷ lệ mỡ, khối lượng cơ, nước, xương, v.v. Dữ liệu được lưu trữ có tổ chức và có thể tái sử dụng cho các mục đích phân tích sau này.
  
  \item \textbf{Phân tích và nội suy thông số sinh học}: Thông qua các công thức y sinh học và kỹ thuật nội suy, hệ thống cung cấp nhiều chỉ số có giá trị đánh giá tình trạng sức khỏe tổng quát của người dùng như BMI, TDEE, BMR, LBM, v.v.
  
  \item \textbf{Ứng dụng trí tuệ nhân tạo (AI)}: Các mô hình học máy được tích hợp vào hệ thống nhằm dự đoán giới tính, tỷ lệ mỡ cơ thể, đo lường độ thăng bằng và cung cấp lời khuyên sức khỏe. So sánh với kết quả từ các thiết bị có cảm biến thực tế cho thấy mức độ chính xác tương đối cao.
  
  \item \textbf{Trực quan hóa và tương tác}: Dữ liệu sau khi được xử lý có thể hiển thị trực tiếp trên nền tảng IoT Core thông qua giao diện bảng điều khiển (Dashboard), giúp người dùng dễ dàng theo dõi sự thay đổi sức khỏe theo thời gian.
  
  \item \textbf{Khả năng mở rộng}: Hệ thống được thiết kế với kiến trúc mở, cho phép tích hợp thêm các cảm biến mới hoặc các dịch vụ y tế số khác trong tương lai.
\end{itemize}

Tuy nhiên, vẫn còn tồn tại một số điểm cần hoàn thiện, đặc biệt là về hiệu năng trong môi trường có nhiều người dùng đồng thời, khả năng bảo mật dữ liệu khi triển khai thực tế, và độ ổn định trong các điều kiện mạng phức tạp.

\subsection{Hướng phát triển}

Nhằm nâng cao chất lượng và khả năng ứng dụng thực tiễn của hệ thống, một số định hướng phát triển trong tương lai được đề xuất như sau:

\begin{itemize}
  \item \textbf{Tối ưu hóa hiệu năng và trải nghiệm người dùng}: Cải tiến tốc độ xử lý dữ liệu và giao diện hiển thị để phù hợp với người dùng cuối, đặc biệt là nhóm người cao tuổi hoặc có trình độ công nghệ hạn chế.
  
  \item \textbf{Bổ sung tính năng chẩn đoán và tư vấn chuyên sâu}: Mở rộng mô hình AI để hỗ trợ nhận diện dấu hiệu bất thường và tự động gợi ý các biện pháp can thiệp ban đầu hoặc đề xuất gặp bác sĩ khi cần thiết.
  
  \item \textbf{Tăng cường bảo mật và bảo vệ quyền riêng tư}: Áp dụng các tiêu chuẩn bảo mật tiên tiến như mã hóa đầu cuối, xác thực hai lớp và kiểm soát truy cập dựa trên vai trò để đảm bảo an toàn dữ liệu người dùng.
  
  \item \textbf{Tích hợp đa nền tảng}: Phát triển ứng dụng tương thích với các hệ điều hành phổ biến (Android, iOS, Web) và mở rộng tích hợp với các hệ thống y tế điện tử hiện hành tại Việt Nam như hồ sơ sức khỏe điện tử cá nhân.
  
  \item \textbf{Hợp tác với các cơ sở y tế và doanh nghiệp}: Tiến hành các nghiên cứu ứng dụng thực tế tại bệnh viện, phòng khám hoặc các dự án y tế cộng đồng nhằm kiểm nghiệm và hoàn thiện hệ thống trước khi triển khai đại trà.
\end{itemize}

Như vậy, có thể khẳng định rằng đề tài không chỉ góp phần nâng cao nhận thức về tầm quan trọng của việc số hóa dữ liệu sức khỏe cá nhân mà còn đặt nền tảng cho các nghiên cứu và phát triển tiếp theo trong lĩnh vực y tế số, hướng đến một hệ sinh thái chăm sóc sức khỏe thông minh, tiện lợi và an toàn hơn cho người dùng.
